\section{The projective ambient family after coefficient extension}

\begin{lemma}[Relative projective space over the modeled affine base]
\label{lem:kn-relative-arithmetic-ambient-base-change}
\lean{EnvelopeProperty.Vanishing.RelativeArithmeticAmbientBaseChange.relative_isPullback,EnvelopeProperty.Vanishing.RelativeArithmeticAmbientBaseChange.twist_isVeryAmple}
\leanok
Let $B$ and $C$ be arbitrary commutative $A$-algebras.  Put
$B_C=C\otimes_A B$ and
$Y=\mathbb P_B^d\times_{\operatorname{Spec}A}\operatorname{Spec}C$.
There is a canonical map $Y\to\operatorname{Spec}B_C$ making $Y$
the base change of $\mathbb P_B^d$ along
$\operatorname{Spec}B_C\to\operatorname{Spec}B$.  The pullback of
$\mathcal O_{\mathbb P_B^d}(1)$ is very ample for this map.  The
composite to $\operatorname{Spec}C$ is the original second projection.
No domain, flatness, or reducedness hypotheses occur.
\end{lemma}
\begin{proof}
The spectrum of the tensor product is the cartesian product of
$\operatorname{Spec}C$ and $\operatorname{Spec}B$ over
$\operatorname{Spec}A$.  The original two projections from $Y$, the
first followed by $\mathbb P_B^d\to\operatorname{Spec}B$, have equal
composites to $\operatorname{Spec}A$.  They therefore induce the
canonical map to $\operatorname{Spec}B_C$.  Its two composites are
exactly those original projections by the universal property of the
cartesian product.  Pasting its proposed square with the tensor
cartesian square gives the defining square of $Y$.  Cancellation of
cartesian squares proves that the proposed square is cartesian.
Finally $\mathcal O(1)$ is very ample on $\mathbb P_B^d$ via its
identity embedding.  The actual projective-space base-change
embedding transports this whole line to a very ample line on $Y$.
\end{proof}
